\section{Isomorphisms of ideal lattices with prescribed determinant}
\label{si-section}

The arithmetic cycle construction requires an isomorphism
between two direct sums of fractional ideals, with its
determinant specified as an element of the number field.
An equality of ideal classes alone would not specify this
determinant.  We prove the exact statement by elementary
two-dimensional operations.

Let \(R\) be a Dedekind domain with fraction field \(F\).
All fractional ideals below are nonzero.  Thus they are
invertible, and their products have unique factorizations
into integral powers of the nonzero prime ideals of \(R\).
For a fractional ideal \(I\), write \(v_P(I)\) for its
exponent at \(P\), and set
\[
 I^{-1}=\{x\in F:xI\subset R\}.
\]
In particular \(II^{-1}=R\).

\begin{lemma}\label{si-two-elements}
For fractional ideals \(C,D\), there are
\[
 p\in C\setminus\{0\},\qquad q\in D,\qquad
 r\in D^{-1},\qquad s\in C^{-1}
\]
such that \(ps-qr=1\).
\end{lemma}
\begin{proof}
Choose any nonzero \(p\in C\).
The ideal \(A=(p)C^{-1}\) is a nonzero integral ideal.
There are only finitely many prime ideals containing it;
denote their set by \(\mathcal P\).
For each \(P\in\mathcal P\), choose
\(q_P\in D\setminus PD\).
Such an element exists: a finite set of generators of
\(D\) has minimum \(P\)-valuation \(v_P(D)\), whereas
every element of \(PD\) has strictly larger valuation.

The Chinese remainder theorem gives elements
\(e_P\in R\) which are \(1\) modulo \(P\) and
\(0\) modulo every other member of \(\mathcal P\).
Put \(q=\sum_{P\in\mathcal P}e_Pq_P\in D\).
Its class in \(D/PD\) is the nonzero class of \(q_P\).
Consequently the integral ideal \((q)D^{-1}\) is not
contained in any prime of \(\mathcal P\).
If \(\mathcal P\) is empty, take \(q=0\), since
then \(A=R\).

It follows that
\[
                 (p)C^{-1}+(q)D^{-1}=R.
\]
Indeed a proper ideal on the left would lie in a maximal
ideal containing \(A\), and every such ideal belongs
to \(\mathcal P\).
Writing \(1=ps+qw\) with
\(s\in C^{-1}\) and \(w\in D^{-1}\), and setting
\(r=-w\), proves the assertion.
\end{proof}

\begin{lemma}\label{si-move}
For any fractional ideals \(I,J,C\), there exists
\(M\in\operatorname{SL}_2(F)\) such that
\[
                       M(I\oplus J)
                            =CI\oplus C^{-1}J.
\]
\end{lemma}
\begin{proof}
Apply Lemma~\ref{si-two-elements} to \(C\) and
\(D=CIJ^{-1}\), and put
\[
       M=\begin{pmatrix}p&q\\r&s\end{pmatrix},
       \qquad
       M^{-1}=\begin{pmatrix}s&-q\\-r&p\end{pmatrix}.
\]
The determinant is exactly \(1\).
The containments
\[
 pI\subset CI,\quad qJ\subset CI,\quad
 rI\subset C^{-1}J,\quad sJ\subset C^{-1}J
\]
show one inclusion in the asserted equality.
For the inverse, use
\[
 sCI\subset I,\quad qC^{-1}J\subset I,\quad
 rCI\subset J,\quad pC^{-1}J\subset J.
\]
These follow immediately from the four ideal memberships
in Lemma~\ref{si-two-elements} and
\(D^{-1}=C^{-1}I^{-1}J\).
Thus the restriction of \(M\) is an isomorphism of the
specified \(R\)-modules.
\end{proof}

\begin{proposition}\label{si-equal-product}
Let \(I_1,\ldots,I_n\) and \(J_1,\ldots,J_n\) be
fractional ideals with
\[
                         \prod_{i=1}^n I_i
                              =\prod_{i=1}^n J_i.
\]
Then an element of \(\operatorname{SL}_n(F)\) maps
\(\bigoplus_i I_i\) onto \(\bigoplus_i J_i\).
For \(n=0\), use the unique automorphism of the zero
space, whose determinant is \(1\).
\end{proposition}
\begin{proof}
For \(n=1\) the ideals are equal and the identity works.
For \(n\ge2\), designate the last coordinate as an
auxiliary coordinate.
For each \(i<n\), let \(I'_i\) be the ideal currently
in coordinate \(i\), and put \(C=J_i(I'_i)^{-1}\).
Apply Lemma~\ref{si-move} in coordinates \(i,n\),
acting as the identity on the other coordinates.
The new ideal in coordinate \(i\) is \(J_i\),
the product of all coordinate ideals is unchanged,
and the determinant of the operation is \(1\).
After these \(n-1\) operations, the first \(n-1\)
coordinate ideals are the required ones.
Cancellation of their invertible product shows that
the last ideal is \(J_n\) as well.
The product of the matrices has determinant \(1\).
\end{proof}

\begin{theorem}\label{si-prescribed}
Let \(I_i,J_i\), \(1\le i\le n\), be fractional
ideals, and let \(\beta\in F^\times\).
There exists \(g\in\operatorname{GL}_n(F)\) with
\[
          g\Bigl(\bigoplus_i I_i\Bigr)
                   =\bigoplus_i J_i,
          \qquad \det g=\beta,
\]
if and only if
\[
             \prod_iJ_i=(\beta)\prod_iI_i,
\]
provided that \(n>0\).  For \(n=0\), existence is
equivalent to \(\beta=1\).
\end{theorem}
\begin{proof}
For necessity when \(n>0\), take the top exterior
power inside \(\bigwedge^n_F F^n=F\).
The exterior power of the first ideal lattice is
\(\prod_iI_i\), that of the second is
\(\prod_iJ_i\), and the induced map is multiplication
by \(\det g\).

For sufficiency, multiply the first coordinate by
\(\beta\).  This has determinant \(\beta\) and
replaces the product of the source ideals by
\((\beta)\prod_iI_i=\prod_iJ_i\).
Apply Proposition~\ref{si-equal-product} afterwards.
The composite has the required image and determinant.
The statement in dimension zero follows from the
definition of the empty determinant.
\end{proof}

To match the source, specialize to
\(R=\mathcal O_F\) and let
\(a_i,b_i\) be finitely supported integral divisors.
Define
\[
 I(a_i)=\prod_P P^{a_i(P)},\qquad
 I(b_i)=\prod_P P^{b_i(P)}.
\]
Membership in \(I(a_i)\) means
\(v_P(x)\ge a_i(P)\) at every prime, with \(x=0\)
always allowed.  These are exactly the valuation
inequalities defining the source's ideal lattice.
The hypothesis
\[
              v_P(\beta)=\sum_i\bigl(b_i(P)-a_i(P)\bigr)
              \qquad\text{for every }P
\]
is exactly the equality of fractional ideals in
Theorem~\ref{si-prescribed}.
Thus the theorem gives the required isomorphism of ideal lattices,
including the condition that the index set be nonempty or \(\beta=1\).
There is no assumption that an individual ideal is
principal.

Both qualifications matter.  In dimension one an
isomorphism of the required kind would force the
ratio ideal to be principal, so arbitrary pairs of
ideals cannot replace the displayed hypothesis.
In dimension zero all divisor sums vanish, so the
divisor hypothesis alone would permit any unit
\(\beta\); nevertheless the determinant remains
\(1\).  For example, \(\beta=-1\) over \(\mathbb Q\)
would contradict the conclusion if the source's
empty-index condition were deleted.
Neither example satisfies all the actual hypotheses.


